\paragrafos{Τετράπλευρα, παραλληλόγραμμα, τραπέζια}{Β.3.10, «Τετράπλευρα, παραλληλόγραμμα, τραπέζια»}

\begin{anagke}
Η θεία Μαρία στρώνει πλακάκια στην αυλή και έχει μπροστά της τρεις στοίβες. Τα πλακάκια της πρώτης έχουν τα απέναντι ζευγάρια πλευρών παράλληλα και εφαρμόζουν τέλεια το ένα δίπλα στο άλλο. Της δεύτερης έχουν μόνο ένα ζευγάρι παράλληλων πλευρών και μοιάζουν με μικρά τραπεζάκια. Της τρίτης δεν έχουν καθόλου παράλληλες πλευρές και δύσκολα «κάθονται» με τα άλλα. «Και τα τρία έχουν τέσσερις πλευρές», λέει η Αγγελική, «αλλά δεν είναι το ίδιο σχήμα». Ο Στέλιος προσπαθεί να τα βάλει όλα στην ίδια κατηγορία, η Ρήνη ρωτά ποια ιδιότητα κάνει τα πλακάκια της πρώτης στοίβας τόσο βολικά, και η Ζωή θέλει να ξέρει γιατί.

Σε αυτή την παράγραφο ονομάζουμε τα τετράπλευρα, μελετάμε το παραλληλόγραμμο και το τραπέζιο και μαθαίνουμε τις ιδιότητές τους.
\end{anagke}

\begin{lexi}{τετράπλευρο, παραλληλόγραμμο, τραπέζιο, διαγώνιος}{ονόματα τετραπλεύρων}
Το \lexiword{τετράπλευρο} έχει τέσσερις πλευρές. Το \lexiword{παραλληλόγραμμο} είναι το σχήμα που ορίζεται από «παράλληλες γραμμές». Το \lexiword{τραπέζιο} έρχεται από το «τράπεζα», το τραπέζι, γιατί μοιάζει με τραπεζάκι. Η \lexiword{διαγώνιος} είναι το τμήμα που περνά «δια γωνίας», από γωνία σε γωνία. Ισοσκελές σημαίνει «με ίσα σκέλη».
\end{lexi}

\begin{metafrasi}
Σχήμα με τέσσερις ευθύγραμμες πλευρές. & τετράπλευρο\\
Τα απέναντι ζεύγη πλευρών είναι παράλληλα. & παραλληλόγραμμο ($\mathrm{ΑΒ}\parallel\mathrm{ΓΔ}$, $\mathrm{ΑΔ}\parallel\mathrm{ΒΓ}$)\\
Μόνο ένα ζεύγος παράλληλων πλευρών. & τραπέζιο\\
Τραπέζιο με ίσες τις μη παράλληλες πλευρές. & ισοσκελές τραπέζιο\\
Η απόσταση ανάμεσα στις παράλληλες πλευρές. & ύψος\\
\end{metafrasi}

\tmima{Η θεωρία}

\begin{orismos}{τετράπλευρο}
\textbf{Τετράπλευρο} λέγεται το ευθύγραμμο σχήμα με τέσσερις πλευρές. Έχει τέσσερις κορυφές και τέσσερις γωνίες. Τα τμήματα που ενώνουν απέναντι κορυφές λέγονται \textbf{διαγώνιοι}.
\end{orismos}

\begin{center}
\begin{tikzpicture}[scale=0.8]
\draw[thick] (0.3,2.2)--(3.6,2.6)--(4,0.2)--(0,0)--cycle;
\draw[gblue,dashed] (0.3,2.2)--(4,0.2); \draw[gblue,dashed] (3.6,2.6)--(0,0);
\node[above left] at (0.3,2.2) {$\mathrm{Α}$}; \node[above right] at (3.6,2.6) {$\mathrm{Β}$};
\node[below right] at (4,0.2) {$\mathrm{Γ}$}; \node[below left] at (0,0) {$\mathrm{Δ}$};
\end{tikzpicture}
\end{center}

\begin{kanonas}{άθροισμα γωνιών τετραπλεύρου}
Μια διαγώνιος χωρίζει το τετράπλευρο σε δύο τρίγωνα, καθένα με άθροισμα γωνιών $180^\circ$. Άρα το άθροισμα των γωνιών κάθε τετραπλεύρου είναι $2\cdot180^\circ=360^\circ$.
\end{kanonas}

\begin{orismos}{παραλληλόγραμμο}
\textbf{Παραλληλόγραμμο} λέγεται το τετράπλευρο που έχει τις απέναντι πλευρές του παράλληλες: $\mathrm{ΑΒ}\parallel\mathrm{ΔΓ}$ και $\mathrm{ΑΔ}\parallel\mathrm{ΒΓ}$.
\end{orismos}

Ειδικές περιπτώσεις παραλληλογράμμου είναι το ορθογώνιο, ο ρόμβος και το τετράγωνο, που θα δούμε στην επόμενη παράγραφο.

\begin{center}
\begin{tikzpicture}[scale=0.9]
\draw[thick] (0,0)--(4,0)--(5.2,2)--(1.2,2)--cycle;
\draw[gblue,dashed] (1.2,2)--(4,0); \draw[gblue,dashed] (5.2,2)--(0,0);
\draw[orange!80!black,thick] (1.2,2)--(1.2,0); \draw (1.2,0.2)--(1.4,0.2)--(1.4,0);
\fill (2.6,1) circle (1.6pt) node[above right] {\small $\mathrm{Ο}$};
\node[above left] at (1.2,2) {$\mathrm{Α}$}; \node[above right] at (5.2,2) {$\mathrm{Β}$};
\node[below right] at (4,0) {$\mathrm{Γ}$}; \node[below left] at (0,0) {$\mathrm{Δ}$};
\node[left,orange!80!black] at (1.2,1) {\small $\upsilon$};
\end{tikzpicture}
\end{center}

\begin{idiotita}{παραλληλογράμμου}
Σε κάθε παραλληλόγραμμο $\mathrm{ΑΒΓΔ}$ ισχύουν:
\begin{itemize}[leftmargin=1.2em,itemsep=0.3ex,topsep=2pt]
\item οι \textbf{απέναντι πλευρές} είναι ίσες: $\mathrm{ΑΒ}=\mathrm{ΓΔ}$ και $\mathrm{ΑΔ}=\mathrm{ΒΓ}$,
\item οι \textbf{απέναντι γωνίες} είναι ίσες: $\hat{\mathrm{Α}}=\hat{\mathrm{Γ}}$ και $\hat{\mathrm{Β}}=\hat{\mathrm{Δ}}$,
\item οι \textbf{διαγώνιοι διχοτομούνται}: $\mathrm{ΟΑ}=\mathrm{ΟΓ}$ και $\mathrm{ΟΒ}=\mathrm{ΟΔ}$,
\item οι \textbf{διαδοχικές γωνίες} είναι παραπληρωματικές: $\hat{\mathrm{Α}}+\hat{\mathrm{Δ}}=180^\circ$.
\end{itemize}
\end{idiotita}

Η τελευταία ιδιότητα προκύπτει από όσα ξέρουμε για τις παράλληλες: οι $\hat{\mathrm{Α}}$ και $\hat{\mathrm{Δ}}$ είναι εντός και επί τα αυτά γωνίες των παραλλήλων $\mathrm{ΑΒ}$, $\mathrm{ΔΓ}$ με τέμνουσα την $\mathrm{ΑΔ}$.

Το σημείο τομής $\mathrm{Ο}$ των διαγωνίων λέγεται \textbf{κέντρο} του παραλληλογράμμου. Κάθε πλευρά μπορεί να ονομαστεί \textbf{βάση}, και η απόσταση της βάσης από την απέναντι πλευρά λέγεται \textbf{ύψος} ($\upsilon$ στο σχήμα). Η \textbf{περίμετρος} είναι το άθροισμα των πλευρών, δηλαδή $2\cdot(\mathrm{ΑΒ}+\mathrm{ΒΓ})$.

\begin{orismos}{τραπέζιο}
\textbf{Τραπέζιο} λέγεται το τετράπλευρο που έχει \emph{μόνο δύο} πλευρές παράλληλες. Οι παράλληλες λέγονται \textbf{βάσεις} (η μεγάλη και η μικρή) και η απόστασή τους \textbf{ύψος}. Αν οι άλλες δύο πλευρές είναι ίσες, το τραπέζιο λέγεται \textbf{ισοσκελές}.
\end{orismos}

\begin{center}
\begin{tikzpicture}[scale=0.75]
\draw[thick] (0,0)--(5,0)--(3.6,2)--(1,2)--cycle;
\node[below left] at (0,0) {$\mathrm{Δ}$}; \node[below right] at (5,0) {$\mathrm{Γ}$};
\node[above right] at (3.6,2) {$\mathrm{Β}$}; \node[above left] at (1,2) {$\mathrm{Α}$};
\node at (2.5,-0.7) {\small τραπέζιο ($\mathrm{ΑΒ}\parallel\mathrm{ΔΓ}$)};
\begin{scope}[xshift=7.5cm]
\draw[thick] (0,0)--(5,0)--(4,2)--(1,2)--cycle;
\draw[dashed,gblue] (2.5,2.5)--(2.5,-0.5);
\fill (2.5,2) circle (1.4pt) node[above right] {\small $\mathrm{Μ}_1$}; \fill (2.5,0) circle (1.4pt) node[below right] {\small $\mathrm{Μ}_2$};
\node[below left] at (0,0) {$\mathrm{Δ}$}; \node[below right] at (5,0) {$\mathrm{Γ}$};
\node[above right] at (4,2) {$\mathrm{Β}$}; \node[above left] at (1,2) {$\mathrm{Α}$};
\node at (2.5,-1.05) {\small ισοσκελές τραπέζιο};
\end{scope}
\end{tikzpicture}
\end{center}

\begin{idiotita}{ισοσκελούς τραπεζίου}
Σε ισοσκελές τραπέζιο $\mathrm{ΑΒΓΔ}$ ($\mathrm{ΑΒ}\parallel\mathrm{ΔΓ}$, $\mathrm{ΑΔ}=\mathrm{ΒΓ}$):
\begin{itemize}[leftmargin=1.2em,itemsep=0.3ex,topsep=2pt]
\item οι γωνίες που βρίσκονται στην ίδια βάση είναι ίσες: $\hat{\mathrm{Δ}}=\hat{\mathrm{Γ}}$ και $\hat{\mathrm{Α}}=\hat{\mathrm{Β}}$,
\item η ευθεία των μέσων $\mathrm{Μ}_1$, $\mathrm{Μ}_2$ των βάσεων είναι \textbf{άξονας συμμετρίας} και μεσοκάθετος των δύο βάσεων.
\end{itemize}
\end{idiotita}

Σε κάθε τραπέζιο οι γωνίες που βρίσκονται στην ίδια μη παράλληλη πλευρά είναι παραπληρωματικές (εντός και επί τα αυτά γωνίες των βάσεων): $\hat{\mathrm{Α}}+\hat{\mathrm{Δ}}=180^\circ$ και $\hat{\mathrm{Β}}+\hat{\mathrm{Γ}}=180^\circ$.

\begin{prosoxi}
\begin{itemize}[leftmargin=1.2em,itemsep=0.3ex,topsep=2pt]
\item Το τραπέζιο έχει \emph{ακριβώς} ένα ζεύγος παράλληλων πλευρών. Αν είχε και το δεύτερο, θα ήταν παραλληλόγραμμο.
\item Το ισοσκελές τραπέζιο έχει ίσες τις μη παράλληλες πλευρές, αλλά δεν είναι παραλληλόγραμμο.
\item Οι διαγώνιοι ενός παραλληλογράμμου διχοτομούνται, αλλά δεν είναι γενικά ίσες ούτε κάθετες.
\item Το ύψος είναι \emph{απόσταση}, κάθετη στις βάσεις, και συνήθως δεν είναι πλευρά του σχήματος.
\end{itemize}
\end{prosoxi}

\tmima{Λυμένα παραδείγματα}

\paradeigma{Πλευρές και περίμετρος}
\tixero Παραλληλόγραμμο $\mathrm{ΑΒΓΔ}$ με $\mathrm{ΑΒ}=7$ cm και $\mathrm{ΒΓ}=4$ cm.\par
\tithelo Οι άλλες δύο πλευρές και η περίμετρος.\par
\possindeo Οι απέναντι πλευρές είναι ίσες: $\mathrm{ΓΔ}=\mathrm{ΑΒ}=7$ cm και $\mathrm{ΑΔ}=\mathrm{ΒΓ}=4$ cm. Η περίμετρος είναι $7+4+7+4=22$ cm.\par
\apantisi $\mathrm{ΓΔ}=7$ cm, $\mathrm{ΑΔ}=4$ cm, περίμετρος $22$ cm.\par
\elegxos $2\cdot(7+4)=22$ cm.

\paradeigma{Γωνίες παραλληλογράμμου}
\tixero Παραλληλόγραμμο $\mathrm{ΑΒΓΔ}$ με $\hat{\mathrm{Α}}=65^\circ$.\par
\tithelo Οι άλλες γωνίες.\par
\possindeo Η $\hat{\mathrm{Γ}}$ είναι απέναντι από την $\hat{\mathrm{Α}}$, άρα $\hat{\mathrm{Γ}}=65^\circ$. Οι $\hat{\mathrm{Β}}$ και $\hat{\mathrm{Δ}}$ είναι διαδοχικές της $\hat{\mathrm{Α}}$, άρα $\hat{\mathrm{Β}}=\hat{\mathrm{Δ}}=180^\circ-65^\circ=115^\circ$.\par
\apantisi $65^\circ,115^\circ,65^\circ,115^\circ$.\par
\elegxos $65+115+65+115=360$.

\paradeigma{Οι διαγώνιοι}
\tixero Στο παραλληλόγραμμο $\mathrm{ΑΒΓΔ}$ οι διαγώνιοι τέμνονται στο $\mathrm{Ο}$ και $\mathrm{ΟΑ}=3{,}5$ cm, $\mathrm{ΟΒ}=2{,}4$ cm.\par
\tithelo Τα μήκη των διαγωνίων.\par
\possindeo Οι διαγώνιοι διχοτομούνται, άρα $\mathrm{ΑΓ}=2\cdot3{,}5=7$ cm και $\mathrm{ΒΔ}=2\cdot2{,}4=4{,}8$ cm.\par
\apantisi $\mathrm{ΑΓ}=7$ cm και $\mathrm{ΒΔ}=4{,}8$ cm.\par
\elegxos Το $\mathrm{Ο}$ είναι μέσο και των δύο διαγωνίων.

\paradeigma{Ισοσκελές τραπέζιο}
\tixero Ισοσκελές τραπέζιο $\mathrm{ΑΒΓΔ}$ με $\mathrm{ΑΒ}\parallel\mathrm{ΔΓ}$ και $\hat{\mathrm{Δ}}=70^\circ$.\par
\tithelo Οι άλλες γωνίες.\par
\possindeo Οι γωνίες της βάσης $\mathrm{ΔΓ}$ είναι ίσες: $\hat{\mathrm{Γ}}=70^\circ$. Οι $\hat{\mathrm{Α}}$ και $\hat{\mathrm{Δ}}$ είναι παραπληρωματικές, άρα $\hat{\mathrm{Α}}=180^\circ-70^\circ=110^\circ$, και $\hat{\mathrm{Β}}=\hat{\mathrm{Α}}=110^\circ$.\par
\apantisi $\hat{\mathrm{Α}}=\hat{\mathrm{Β}}=110^\circ$ και $\hat{\mathrm{Γ}}=\hat{\mathrm{Δ}}=70^\circ$.\par
\elegxos $110+110+70+70=360$.

\paradeigma{Τι είδους τετράπλευρο;}
\tixero Στο τετράπλευρο $\mathrm{ΑΒΓΔ}$ είναι $\hat{\mathrm{Α}}=105^\circ$, $\hat{\mathrm{Β}}=100^\circ$ και $\hat{\mathrm{Δ}}=75^\circ$.\par
\tithelo Η $\hat{\mathrm{Γ}}$ και το είδος του τετραπλεύρου.\par
\possindeo Το άθροισμα είναι $360^\circ$, άρα $\hat{\mathrm{Γ}}=360^\circ-(105^\circ+100^\circ+75^\circ)=80^\circ$. Με τέμνουσα την $\mathrm{ΑΔ}$ οι $\hat{\mathrm{Α}}$, $\hat{\mathrm{Δ}}$ έχουν άθροισμα $180^\circ$, άρα $\mathrm{ΑΒ}\parallel\mathrm{ΔΓ}$. Με τέμνουσα την $\mathrm{ΑΒ}$ οι $\hat{\mathrm{Α}}$, $\hat{\mathrm{Β}}$ έχουν άθροισμα $205^\circ\neq180^\circ$, άρα οι $\mathrm{ΑΔ}$ και $\mathrm{ΒΓ}$ δεν είναι παράλληλες.\par
\apantisi $\hat{\mathrm{Γ}}=80^\circ$ και το τετράπλευρο είναι τραπέζιο (όχι παραλληλόγραμμο).\par
\elegxos Στο παραλληλόγραμμο θα ήταν $\hat{\mathrm{Α}}+\hat{\mathrm{Β}}=180^\circ$.

\tmima{Ασκήσεις}

\zoni{Το κατάλαβα}{ονόματα και ιδιότητες}

\ex Σωστό ή λάθος; α) Κάθε παραλληλόγραμμο είναι τετράπλευρο. β) Κάθε τετράπλευρο είναι παραλληλόγραμμο. γ) Το τραπέζιο έχει ακριβώς ένα ζεύγος παράλληλων πλευρών. δ) Οι διαγώνιοι του παραλληλογράμμου διχοτομούνται. ε) Οι διαγώνιοι κάθε παραλληλογράμμου είναι ίσες. στ) Οι απέναντι γωνίες παραλληλογράμμου είναι ίσες.

\ex Σχεδίασε ένα παραλληλόγραμμο $\mathrm{ΑΒΓΔ}$, τις διαγωνίους του και το κέντρο $\mathrm{Ο}$. Μέτρησε τα $\mathrm{ΟΑ}$ και $\mathrm{ΟΓ}$, και φέρε ένα ύψος του.

\ex Στο παραλληλόγραμμο $\mathrm{ΑΒΓΔ}$ είναι $\mathrm{ΑΒ}=9$ cm και $\mathrm{ΒΓ}=5$ cm. Βρες τις $\mathrm{ΓΔ}$, $\mathrm{ΑΔ}$ και την περίμετρο.

\ex Στο παραλληλόγραμμο $\mathrm{ΑΒΓΔ}$ είναι $\hat{\mathrm{Α}}=72^\circ$. Βρες τις άλλες γωνίες.

\ex Οι διαγώνιοι παραλληλογράμμου $\mathrm{ΑΒΓΔ}$ τέμνονται στο $\mathrm{Ο}$ και είναι $\mathrm{ΑΓ}=10$ cm, $\mathrm{ΒΔ}=8$ cm. Βρες τα $\mathrm{ΟΑ}$, $\mathrm{ΟΓ}$, $\mathrm{ΟΒ}$, $\mathrm{ΟΔ}$.

\ex Στο ισοσκελές τραπέζιο $\mathrm{ΑΒΓΔ}$ ($\mathrm{ΑΒ}\parallel\mathrm{ΔΓ}$) είναι $\hat{\mathrm{Δ}}=68^\circ$. Βρες τις άλλες γωνίες.

\ex Στο τραπέζιο $\mathrm{ΑΒΓΔ}$ ($\mathrm{ΑΒ}\parallel\mathrm{ΔΓ}$) είναι $\hat{\mathrm{Α}}=95^\circ$ και $\hat{\mathrm{Β}}=110^\circ$. Βρες τις $\hat{\mathrm{Δ}}$ και $\hat{\mathrm{Γ}}$.

\zoni{Το χρησιμοποιώ}{πλακάκια και ζαρντινιέρες}

\ex \textbf{Το πλακάκι.} Ένα πλακάκι της θείας Μαρίας είναι παραλληλόγραμμο με πλευρές $30$ cm και $20$ cm. Θέλει να κολλήσει ταινία γύρω από την άκρη του. Πόση ταινία χρειάζεται;

\ex \textbf{Η ζαρντινιέρα.} Η πλαϊνή όψη μιας ζαρντινιέρας είναι ισοσκελές τραπέζιο με μικρή βάση $40$ cm, μεγάλη βάση $60$ cm και μη παράλληλες πλευρές $25$ cm. Οι γωνίες στη μεγάλη βάση είναι $70^\circ$. Βρες τις άλλες γωνίες και την περίμετρο.

\ex \textbf{Ο μηχανισμός του Ηλία.} Ο Ηλίας έχει τέσσερις ράβδους που ενώνονται με καρφιά σε παραλληλόγραμμο. Στην αρχή η $\hat{\mathrm{Α}}$ είναι $90^\circ$. Τον πιέζει ώστε να γίνει $\hat{\mathrm{Α}}=60^\circ$. Βρες τις άλλες γωνίες. Αλλάζουν τα μήκη των πλευρών;

\ex \textbf{Η περίμετρος.} Ένα παραλληλόγραμμο οικόπεδο έχει περίμετρο $46$ m και η μια πλευρά του είναι κατά $3$ m μεγαλύτερη από την άλλη. Βρες τις πλευρές του.

\ex \textbf{Το ύψος.} Σχεδίασε παραλληλόγραμμο $\mathrm{ΑΒΓΔ}$ με $\mathrm{ΑΒ}=5$ cm, $\mathrm{ΑΔ}=3$ cm και $\hat{\mathrm{Α}}=60^\circ$. Φέρε το ύψος από το $\mathrm{Δ}$ προς την $\mathrm{ΑΒ}$ και μέτρησέ το.

\zoni{Το σκέφτομαι}{γιατί και πώς}

\ex Εξήγησε γιατί σε κάθε παραλληλόγραμμο $\mathrm{ΑΒΓΔ}$ οι γωνίες $\hat{\mathrm{Α}}$ και $\hat{\mathrm{Δ}}$ έχουν άθροισμα $180^\circ$.

\ex Το τετράπλευρο $\mathrm{ΑΒΓΔ}$ έχει $\hat{\mathrm{Α}}=110^\circ$ και $\hat{\mathrm{Δ}}=65^\circ$. Μπορεί να είναι παραλληλόγραμμο; Δικαιολόγησε.

\ex Στο τετράπλευρο $\mathrm{ΕΖΗΔ}$ είναι $\hat{\mathrm{Δ}}=\hat{\mathrm{Η}}=90^\circ$, $\mathrm{ΔΕ}=4$ cm και $\mathrm{ΗΖ}=7$ cm. Δείξε ότι $\mathrm{ΔΕ}\parallel\mathrm{ΗΖ}$ και ότι το τετράπλευρο είναι τραπέζιο και όχι παραλληλόγραμμο.

\ex Στο ισοσκελές τραπέζιο $\mathrm{ΑΒΓΔ}$ ($\mathrm{ΑΒ}\parallel\mathrm{ΔΓ}$) είναι $\hat{\mathrm{Δ}}=\hat{\mathrm{Γ}}$. Εξήγησε γιατί $\hat{\mathrm{Α}}=\hat{\mathrm{Β}}$.

\zoni{Για το σπίτι}{ρίξε μια ματιά}

\ex Βρες στο σπίτι ή στον δρόμο τρία αντικείμενα που έχουν σχήμα παραλληλογράμμου ή τραπεζίου (παράθυρο, πλακάκι, ράφι, κερκίδα). Σχεδίασέ τα και σημείωσε ποιες πλευρές είναι παράλληλες.

\ex Σχεδίασε παραλληλόγραμμο $\mathrm{ΑΒΓΔ}$ με $\mathrm{ΑΒ}=6$ cm, $\mathrm{ΑΔ}=4$ cm και $\hat{\mathrm{Α}}=70^\circ$. Βρες με υπολογισμό και με μέτρηση τις άλλες γωνίες του.
